\section{Introduction}
\label{sec:intro}

\begin{sourceblock}{Opening Paragraph - Traditional ML Assumptions}
Traditional machine learning paradigms operate under the fundamental assumption that data features are passively observed and completely reported. In many socio-economic applications, however, the data is generated by strategic agents who possess private preferences and can actively manipulate the information they share. When a decision-maker commits to a predictive model, these agents strategically decide which features to reveal and which to withhold to steer the model's prediction in their favor. This phenomenon of voluntary disclosure under strategic feature withholding is common in markets characterized by asymmetric information, where the decision-maker must predict an outcome based on partially self-reported feature vectors.
\end{sourceblock}

\begin{faqblock}
\faqitem{Context}{What assumption does traditional ML make about data?}
\faqitem{Causes}{Why do strategic agents manipulate information?}
\faqitem{Consequence}{What is voluntary disclosure?}
\faqanswers
\faqans{Context}{That features are passively observed and completely reported.}
\faqans{Causes}{They have private preferences and want to steer model predictions.}
\faqans{Consequence}{Strategic agents decide which features to reveal/withhold to optimize their outcomes.}
\end{faqblock}

\begin{sourceblock}{Insurance Example - Multi-directional Incentives}
To illustrate the multi-directional incentives of such strategic withholding, consider an individual navigating the insurance sector. When purchasing a health or life insurance policy, the applicant seeks to minimize their premium. They have a strategic incentive to withhold negative health indicators---such as chronic conditions or high-risk habits---to present an understated risk profile to the underwriter. Conversely, if the same individual later files a claim for disability benefits or worker's compensation, their preference flips: they now seek to maximize the payout, motivating them to selectively disclose their medical history to present an overstated severity of their condition. In both scenarios, the observed features are not missing at random; they are the output of a utility-maximizing disclosure strategy that directly responds to the committed prediction rule.
\end{sourceblock}

\begin{faqblock}
\faqitem{Context}{What incentives exist when buying insurance?}
\faqitem{Process}{How do incentives change when filing claims?}
\faqitem{Consequence}{What determines observed features?}
\faqanswers
\faqans{Context}{Minimize premium by withholding negative health indicators.}
\faqans{Process}{Flip to maximizing payout by revealing/exaggerating severity.}
\faqans{Consequence}{Utility-maximizing disclosure strategy responding to the prediction rule.}
\end{faqblock}

\begin{sourceblock}{Research Scope and Questions}
Crucially, strategic feature withholding is not universally detrimental or beneficial to the system. Whether voluntary disclosure improves or degrades the decision-maker's prediction accuracy depends heavily on the alignment between the sender's preferences and the true target, the inherent value of the features, and external channel noise. We are interested in formally studying and quantifying these dynamics. This motivates a dual-level investigation: first, a \textit{decision-theoretic} game analysis to characterize when voluntary disclosure yields a positive strategic advantage over mandatory disclosure; and second, a \textit{learning-theoretic} analysis to determine how the decision-maker can learn the optimal predictive rules under finite sample regimes (offline) and adapt dynamically in real-time under partial or bandit feedback (online).
\end{sourceblock}

\begin{faqblock}
\faqitem{Context}{Does voluntary disclosure always help or hurt?}
\faqitem{Process}{What factors determine its impact?}
\faqitem{Conclusion}{What are the two levels of investigation?}
\faqanswers
\faqans{Context}{Neither; impact depends on alignment, feature value, and noise.}
\faqans{Process}{Sender-target alignment, feature informativeness, channel noise levels.}
\faqans{Conclusion}{(1) Decision-theoretic: when does voluntary disclosure help? (2) Learning-theoretic: how to learn optimal rules offline and online.}
\end{faqblock}

\begin{sourceblock}{Main Contributions - Regression Results}
In this paper, we establish a comprehensive set of theoretical results across the six settings of strategic feature withholding, outlining the fundamental limits, sample complexities, and online regret behaviors. For the regression setting, we establish necessary and sufficient conditions for the decision-theoretic game under the structural assumption of an all-or-nothing disclosure setup where agents reveal the entire feature vector or withhold it completely. We show that when the sender's preferences are aligned with the true target, the strategic advantage is positive because the sender naturally reveals features when they are most predictive. In the offline batch learning regime, we prove that the strategic empirical risk minimizer (ERM) on a smooth surrogate loss is consistent and achieves a uniform generalization rate of $\mathcal{O}(d^3 / \sqrt{n})$ under finite training samples. For the online setting, we design a gradient-free online learning algorithm that operates under bandit feedback and achieves a cumulative regret of $\mathcal{O}(d \sqrt{T})$, establishing that the receiver can dynamically learn the strategic prediction rule in real-time despite the partial observability of features.
\end{sourceblock}

\begin{faqblock}
\faqitem{Context}{What are the three regimes analyzed for regression?}
\faqitem{Consequence}{What is the strategic advantage result for regression?}
\faqitem{Consequence}{What are the learning rates for regression?}
\faqanswers
\faqans{Context}{(1) Decision-theoretic game; (2) Offline batch learning; (3) Online bandit learning.}
\faqans{Consequence}{Strategic advantage is positive when sender preferences align with true target.}
\faqans{Consequence}{Generalization rate $\mathcal{O}(d^3/\sqrt{n})$; online regret $\mathcal{O}(d\sqrt{T})$.}
\end{faqblock}

\begin{sourceblock}{Main Contributions - Classification Results}
For the multiclass classification setting, we show that the strategic risk is an affine function of the channel noise, where the optimal default class partitions the noise space into a piecewise-linear envelope of at most $K$ subintervals. Under the same all-or-nothing disclosure assumption, we formulate a joint offline learning algorithm and prove that the excess risk decays at a rate of $\mathcal{O}(\sqrt{Kd \log n / n})$ based on the Natarajan dimension of the hypothesis space. In the online sequential classification game, we propose a decoupled EXP3+OGD algorithm that explores the discrete default classes in a bandit fashion while optimizing the continuous predictors, proving an expected cumulative regret of $\tilde{\mathcal{O}}(T^{2/3})$ under a non-degenerate mass condition on the revealed subpopulation.
\end{sourceblock}

\begin{faqblock}
\faqitem{Context}{What structural property exists for classification?}
\faqitem{Consequence}{What is the offline learning rate for classification?}
\faqitem{Consequence}{What is the online regret for classification?}
\faqanswers
\faqans{Context}{Strategic risk is affine in noise; optimal default creates piecewise-linear envelope.}
\faqans{Consequence}{Excess risk decays at $\mathcal{O}(\sqrt{Kd\log n/n})$.}
\faqans{Consequence}{Expected regret $\tilde{\mathcal{O}}(T^{2/3})$ under non-degenerate mass condition.}
\end{faqblock}

\subsection{Related Work}

\begin{sourceblock}{Economic Foundations Overview}
Our work lies at the confluence of microeconomic information design, strategic classification, and learning theory. Rather than analyzing these topics in isolation, we position our work by tracing how our unified framework builds upon and extends three major thematic areas in the literature.
\end{sourceblock}

\begin{faqblock}
\faqitem{Context}{What three areas does this work bridge?}
\faqanswers
\faqans{Context}{Microeconomic information design; strategic classification; learning theory.}
\end{faqblock}

\begin{sourceblock}{Voluntary Disclosure Literature}
First, our work is related to the literature on voluntary disclosure and information design. The economic foundations of voluntary disclosure under verifiable information were pioneered by Milgrom~\cite{milgrom1981good} and Grossman~\cite{grossman1981informational}, who established the classical ``unraveling'' result: when information can be verified, a rational receiver assumes the worst about any withheld information, forcing the sender to disclose all private details in equilibrium. However, complete unraveling fails under realistic market frictions. Dye~\cite{dye1985disclosure} and Verrecchia~\cite{verrecchia1983discretionary} showed that if there is uncertainty regarding whether the sender possesses the information (often modeled as an erasure channel) or if disclosure incurs proprietary costs, senders will strategically withhold unfavorable information. These models were extended to asset pricing by Jovanovic~\cite{jovanovic1982truthful} and Shin~\cite{shin2003disclosures}, and generalized under the framework of Bayesian persuasion by Kamenica and Gentzkow~\cite{kamenica2011bayesian}. While this literature focuses on low-dimensional signals with fixed sender-receiver actions, our work extends these voluntary disclosure models to high-dimensional feature spaces ($X \in \mathbb{R}^d$) and analyzes how a predictive model can be learned. We establish the necessary and sufficient conditions under which voluntary feature disclosure provides a strategic advantage to the receiver, extending Dye's classical results to multidimensional regression and multiclass classification settings.
\end{sourceblock}

\begin{faqblock}
\faqitem{Context}{What is the ``unraveling'' result?}
\faqitem{Causes}{Why does unraveling fail in practice?}
\faqitem{Conclusion}{How does this paper extend prior work?}
\faqanswers
\faqans{Context}{When information is verifiable, rational receivers assume worst about withheld info, forcing full disclosure.}
\faqans{Causes}{Uncertainty about sender's info possession; erasure channels; disclosure costs.}
\faqans{Conclusion}{Extends to high-dimensional feature spaces; establishes conditions for strategic advantage.}
\end{faqblock}

\begin{sourceblock}{Strategic Classification Literature}
Second, our work connects to strategic classification and performative prediction. Within computer science, the interaction between strategic agents and predictive models is studied under strategic classification, pioneered by Hardt et al.~\cite{hardt2016strategic}. In their Stackelberg formulation, agents pay a cost to continuously alter their features to obtain a favorable classification. This has been extended to study algorithmic recourse by Tsirtsis and Rodriguez~\cite{tsirtsis2020decisions}, causal strategic models by Shavit et al.~\cite{shavit2020causal}, and performative prediction by Perdomo et al.~\cite{perdomo2020performative} and Miller et al.~\cite{miller2020strategic}, which analyze how predictions shift the underlying distribution. Further, researchers have investigated the social costs~\cite{milli2019social} and fairness implications~\cite{hu2019disparate} of such models. In contrast to standard strategic classification—where agents continuously perturb features at a cost—our framework focuses on a discrete action space where agents strategically withhold or reveal their entire verifiable feature vector. This discrete choice introduces non-convexities and step-like strategic risks. We show how these discrete interactions partition the noise space $[0,1]$ into piecewise-linear optimal default regions and analyze how the receiver can learn under these discontinuous risk landscapes.
\end{sourceblock}

\begin{faqblock}
\faqitem{Context}{What is standard strategic classification?}
\faqitem{Causes}{How does this paper differ?}
\faqitem{Consequence}{What mathematical challenges arise?}
\faqanswers
\faqans{Context}{Agents pay costs to continuously alter features for favorable classification.}
\faqans{Causes}{Discrete action space: agents withhold or reveal entire feature vectors.}
\faqans{Consequence}{Non-convexities, step-like risks, piecewise-linear default regions.}
\end{faqblock}

\begin{sourceblock}{Learning under Incentive Constraints}
Third, our work is positioned within the literature on learning and optimization under strategic constraints. Dekel et al.~\cite{dekel2010incentive} and Chen et al.~\cite{chen2018incentive} studied offline incentive-compatible regression where agents strategically manipulate their target labels or features. In online settings, Chen et al.~\cite{chen2020learning} and Dong et al.~\cite{dong2018strategic} analyzed online gradient updates when learning from strategic agents. Under bandit feedback or partial observability, online convex optimization relies on gradient-free techniques, such as the one-point bandit gradient estimator of Flaxman et al.~\cite{flaxman2005online}, and multi-armed bandit dynamics as reviewed in Auer et al.~\cite{auer2002nonadversarial} and Cesa-Bianchi and Lugosi~\cite{cesabianchi2006prediction}. We bridge these domains by deriving the first uniform generalization bounds for strategic feature withholding in both offline regression (via product-class Rademacher complexity) and multiclass classification (via Natarajan dimension). In online settings under bandit feedback—where features are hidden when agents withhold—we design updates that combine Flaxman's bandit gradient descent for regression and a decoupled EXP3+OGD algorithm for classification, proving sublinear cumulative regret bounds.
\end{sourceblock}

\begin{faqblock}
\faqitem{Context}{What prior work exists on learning with strategic agents?}
\faqitem{Process}{What gradient-free techniques are used for bandit feedback?}
\faqitem{Conclusion}{What are this paper's key contributions?}
\faqanswers
\faqans{Context}{Incentive-compatible regression; online gradient updates with strategic agents.}
\faqans{Process}{One-point bandit gradient estimators; multi-armed bandit dynamics (EXP3).}
\faqans{Conclusion}{First uniform generalization bounds for strategic withholding; sublinear regret bounds under bandit feedback.}
\end{faqblock}
